\section{Local stability under a centred twist}
\label{sec:r7-twist}

\begin{theorem}[Hyperbolically amplified twist Lipschitz bound]
\label{thm:r7-twist}
For the observation projector $P_L$ and all $\theta,\theta'$ in the reduced
angle interval,
\begin{equation}
 \|P_L[H(\theta)-H(\theta')]P_L\|
 \le C_L|\theta-\theta'|,
 \qquad
 C_L=\mu R\sinh(L/R)C_{\rm Schur}.
 \label{eq:r7-local-lip}
\end{equation}
In particular
\begin{equation}
 C_L\le \frac{\mu R C_{\rm Schur}}2e^{L/R},\qquad
 |\theta-\theta'|\le
 \frac{2\varepsilon}{\mu R C_{\rm Schur}}e^{-L/R}
 \Longrightarrow
 \|P_L(H_\theta-H_{\theta'})P_L\|\le\varepsilon.
 \label{eq:r7-angle-scale}
\end{equation}
\end{theorem}

\begin{proof}
The angular orbit of a point $y$ at radius $\rho$ has speed
\begin{equation*}
 R\sinh(\rho/R).
\end{equation*}
Distance is one-Lipschitz in either
endpoint, and $d\mapsto\sqrt{h^2+d^2}$ has derivative at most one.  Hence,
for $x,y\in P_L$ and every intermediate angle,
\begin{equation}
 |\partial_\theta T_\theta(x,y)|
 \le\mu R\sinh(L/R)|T_\theta(x,y)|.
\end{equation}
The mean-value theorem followed by the uniform Schur estimate proves
\eqref{eq:r7-local-lip}.  The second statement uses
$\sinh u\le e^u/2$.
\end{proof}

Thus the relation $|\delta\theta|e^{L/R}\ll1$ is a consequence of the
microscopic kernel, not a phenomenological moir\'e prescription.

For a finite parameter chart $p$ let
\begin{equation}
 C_p=8|\partial_p t|+
 \sup_\theta\sup_x\sum_y|\partial_pT_\theta(x,y)|.
 \label{eq:r7-Cp}
\end{equation}
Whenever the displayed derivatives exist and obey the same exponential
majorant, $C_p<\infty$ and
$\|P_L[H(\theta,p)-H(\theta,p')]P_L\|
\le C_p\|p-p'\|$.  For example,
$C_w=e^{\mu h}S(\mu)$ and $C_h\le\mu C_{\rm Schur}$; a decay-rate derivative
is bounded by the differentiated geometric series
\begin{equation}
 C_\mu\le |w|e^{\mu h}A_\Gamma a_0
 \frac{e^{-\gamma a_0}}{(1-e^{-\gamma a_0})^2}
 +hC_{\rm Schur}.
\end{equation}
